\section{Models and headline results}
\label{sec:models}

\subsection{Model zoo}

Four families span linear, tree-based and graph-neural backbones, which lets us isolate where the non-linear and graph-structural signal actually matters. Each is trained twice -- once on the 37 pre-departure features (scenario~A), once on the 40 inflight features (scenario~B) -- on identical splits.

\begin{itemize}[leftmargin=*]
    \item \textbf{Logistic regression}~\cite{pedregosa2011scikit} -- L2-regularised binary classifier, SAGA solver, balanced \texttt{class\_weight}, maximum 2{,}000 iterations. The SAGA solver supports both L1 and L2 penalties on million-row tabular data and converges reliably under heavy class imbalance. Strict linear baseline.
    \item \textbf{LightGBM}~\cite{ke2017lightgbm} -- gradient-boosted decision trees with histogram-based, leaf-wise growth: 1{,}500 trees, 63 leaves, depth 7, \texttt{scale\_pos\_weight} = train\_neg/train\_pos $\approx 9.7$, learning rate \num{0.05}, early stopping on val PR-AUC after 50 stagnant rounds.
    \item \textbf{XGBoost}~\cite{chen2016xgboost} -- depth-wise growth on the same hyperparameters: 1{,}500 trees, \texttt{tree\_method='hist'}, \texttt{eval\_metric='aucpr'}, $\eta=\num{0.05}$, $\lambda=1$ (L2), $\alpha=0$ (L1), early stop after 50 stagnant rounds.
    \item \textbf{GraphSAGE}~\cite{hamilton2017inductive} -- a two-layer inductive GNN over the station adjacency graph (\num{2398} nodes, \num{163902} directed edges). Each layer is a \texttt{SAGEConv} with mean aggregation followed by LayerNorm and a ReLU; a Jumping-Knowledge \texttt{cat} block concatenates layer-1 and layer-2 outputs into the per-stop MLP head (32~$\to$~16~$\to$~1, dropout~\num{0.3}). Trained with AdamW (\texttt{lr}~$=\num{1e-3}$, \texttt{weight\_decay}~$=\num{1e-4}$) and a cosine-decay schedule to \num{1e-5} over 50 epochs. PyTorch~\cite{paszke2019pytorch} + PyTorch Geometric~\cite{fey2019fast}.
\end{itemize}

\subsection{Training procedure}
\label{sec:training}

The same loop runs for every (model, scenario) pair. The deliberate symmetry across architectures lets us interpret performance gaps as architectural rather than tuning artefacts.

\begin{enumerate}[leftmargin=*]
    \item \textbf{Loss.} All four optimise weighted binary cross-entropy on the per-stop disruption label. The positive-class weight is the train-set negative-to-positive ratio, $w_+ \approx 9.7$. Minimising raw BCE on the unweighted labels would let any of the four models collapse to ``never disrupted'' and still record \SI{91}{\percent} accuracy; reweighting makes a missed positive 9.7$\times$ more expensive than a missed negative.
    \item \textbf{Optimiser and regularisation.} Logistic regression uses SAGA (a variance-reduced stochastic gradient method that handles non-smooth penalties). LightGBM and XGBoost use second-order gradient-boosted CART; regularisation is implicit in the leaves/depth caps and explicit through XGBoost's $\lambda=1$ L2 leaf-weight penalty. GraphSAGE uses AdamW (decoupled weight decay, which behaves better than L2-on-Adam when LayerNorm scales matter), \texttt{weight\_decay}~$=\num{1e-4}$, dropout~\num{0.3} in the MLP head. None of the four uses early-layer L1 -- the L2 / leaf caps are sufficient given the modest feature count.
    \item \textbf{Learning-rate schedule.} Tabular models are step-wise constant ($\eta=\num{0.05}$ with early stop). GraphSAGE uses a cosine decay from \num{1e-3} to \num{1e-5} over 50 epochs; the long tail at low LR helps the LayerNorm scales settle.
    \item \textbf{Early stopping watches val PR-AUC, not loss.} BCE loss and PR-AUC are correlated but not co-monotonic under heavy imbalance; PR-AUC is the metric that matters operationally, so it is the metric the stopping criterion watches. Stop after 50 stagnant rounds (tabular) or 10 stagnant epochs (GraphSAGE).
    \item \textbf{Threshold tuning.} A static \num{0.5} threshold is rarely optimal under \SI{9}{\percent} prevalence. Each model inherits a \texttt{StopModel.tune\_threshold} method that sweeps 19 candidate thresholds in $[\num{0.05}, \num{0.95}]$ on \emph{validation} probabilities and persists the val-F1-optimal one.
    \item \textbf{Isotonic calibration when needed.} If val ECE $> \num{0.05}$, an isotonic regression~\cite{zadrozny2001calibration} is fit on val probabilities; the threshold is then refit on the calibrated scores. Logistic regression scenario~A is the only model that retains an ECE above \num{0.05} after calibration -- the linear capacity is too low to match the empirical rate at high probabilities.
    \item \textbf{Test set is touched once.} All hyperparameter choices, the early-stop decision, the threshold sweep and the calibration decision use validation only. The test set is read once, after the model is frozen, and produces the headline metrics in \cref{tab:headline}.
\end{enumerate}

\subsection{Wall-clock fit times}

Tabular models train on a multi-core CPU host; GraphSAGE on a single Tesla~T4. Wall-clock fit time per (model, scenario): logreg \SI{123}{s} / \SI{83}{s}, LightGBM \SI{29}{s} / \SI{240}{s}, XGBoost \SI{174}{s} / \SI{1166}{s}, GraphSAGE \SI{74}{s} / \SI{134}{s}. The XGBoost scenario-B blow-up to roughly nineteen minutes is interpretable: with the inflight features the model keeps finding good splits, so the tree count grows substantially before early stopping fires.

\subsection{Test-set headline metrics}

\Cref{tab:headline} reports PR-AUC and F1 on the \num{2706443}-row June~2024 test set; the random-classifier PR-AUC baseline equals the prevalence (\num{0.087}). XGBoost wins both scenarios; the inflight uplift exceeds \num{0.4}~PR-AUC for every architecture.

\begin{table}[H]
    \centering
    \caption{Test-set PR-AUC and F1. ``Uplift'' = PR-AUC$_B$ -- PR-AUC$_A$.}
    \label{tab:headline}
    \footnotesize
    \begin{tabular}{l
        S[table-format=1.3] S[table-format=1.3]
        S[table-format=1.3] S[table-format=1.3]
        S[table-format=1.3]}
        \toprule
        & \multicolumn{2}{c}{Scenario A (pre-departure)} & \multicolumn{2}{c}{Scenario B (inflight)} & \\
        \cmidrule(lr){2-3}\cmidrule(lr){4-5}
        Model & {PR-AUC} & {F1} & {PR-AUC} & {F1} & {Uplift} \\
        \midrule
        Logistic Regression  & 0.481 & 0.458 & 0.881 & 0.820 & 0.400 \\
        LightGBM             & 0.503 & 0.506 & 0.934 & 0.877 & 0.431 \\
        \textbf{XGBoost}     & {\bfseries 0.527} & {\bfseries 0.512} & {\bfseries 0.939} & {\bfseries 0.882} & 0.412 \\
        GraphSAGE            & 0.520 & 0.518 & 0.921 & 0.891 & 0.401 \\
        \bottomrule
    \end{tabular}
\end{table}

\begin{figure}[H]
    \centering
    \includegraphics[width=0.78\textwidth]{stop_level/metrics_summary_B.png}
    \caption{Per-model PR-AUC, F1, precision and recall on the test set under scenario~B. Tree- and graph-based models cluster above \num{0.92}; logreg saturates lower.}
    \label{fig:metrics-summary}
\end{figure}

\paragraph{Calibration.} Post-isotonic test ECE is below \num{0.002} for all four tree models (LightGBM~A \num{0.001}, LightGBM~B \num{0.000}, XGBoost~A \num{0.002}, XGBoost~B \num{0.000}). Logistic regression scenario~A retains ECE \num{0.076}: a residual the linear model's lower information capacity cannot smooth out.

\subsection{Per-country and cascading-respect breakdowns}

\begin{table}[H]
    \centering
    \caption{XGBoost / scenario B, test PR-AUC and F1 by country. Finland advantage: short, recurrent runs whose \texttt{train\_station\_lag7\_rate} is highly predictive.}
    \label{tab:per-country}
    \footnotesize
    \begin{tabular}{l S[table-format=7.0,group-separator={,}] S[table-format=1.3] S[table-format=1.3]}
        \toprule
        Country & {n\_test} & {PR-AUC} & {F1} \\
        \midrule
        Finland     &  495040 & 0.990 & 0.961 \\
        Italy       &  421971 & 0.871 & 0.790 \\
        Netherlands & 1789432 & 0.857 & 0.805 \\
        \bottomrule
    \end{tabular}
\end{table}

\paragraph{Cascading respect (B/XGB).} Conditioning on the previous stop's actual disruption label,
$P(\hat{y}=1 \mid y_{k-1}=1)=\num{0.938}$ versus $P(\hat{y}=1 \mid y_{k-1}=0)=\num{0.005}$, gap $=\num{0.933}$. The model uses the inflight signal aggressively and correctly.

\paragraph{Operating points.} XGBoost~B is operationally usable as a high-precision triage signal: precision at recall \num{0.70} reaches \num{0.985}; at recall \num{0.95} it is still \num{0.548} (compared to \num{0.179} for XGBoost~A at the same recall). Performance is flat across route-position quartiles -- per-quartile PR-AUC stays in $[\num{0.935}, \num{0.954}]$.
